% Abstract
\section*{Abstract}

Concept erasure should remove a target concept without damaging the rest of the representation, but nonlinear concepts remain recoverable after standard linear interventions and unconstrained ambient-space edits can distort representations. We propose \TanGCE{}, a new manifold-aware erasure method that estimates local tangent spaces with neighborhood PCA and restricts erasure updates to that local support. The core idea is simple: if natural representations lie on a low-dimensional manifold inside the ambient space, then effective interventions should follow that manifold rather than push representations away from it. Across our completed text and vision benchmarks, \TanGCE{} delivers the strongest overall nonlinear erasure results in the study: it consistently beats the prior learned baselines, including \AmbGCE{}, and avoids the severe control damage incurred by distortion-based baselines. More broadly, these results suggest that manifold-aware editing is a useful principle not only for concept removal, but for representation editing and optimization in high-dimensional models.
